\scanpage{168}
\mspage{117}
% archive sheet number: 168
n'est autre que $f_!(e_X)$. D'autre part, pour un préfaisceau quelconque \inserted{$G$} sur $I$, on a

\[
g_!(G)\simeq\pi_0(I_{/\unsure{C|G}})\simeq\underset{I}{\varinjlim}\,G(i),
\]
et appliquant ceci à $f_!(e_X)$, on trouve la formule \unsure{($\alpha$)} de la prop.~12.

Revenons à la situation \unsure{avec|où} $f$ $W$-\uline{propre}, et considérons un foncteur

\[
p:X\longrightarrow Y,
\]
nous sommes intéressés à \lq\lq calculer\rq\rq

\[
p_!(e_X)=(y\mapsto\pi_0({}_y\backslash X))\in\mathrm{Ob}\,Y^\wedge .
\]

Considérons pour ceci le diagramme suivant

\diagram{\begin{tikzcd}[ampersand replacement=\&]
{}_y\backslash X_i \arrow[r]\arrow[d] \& {}_y\backslash X \arrow[r]\arrow[d] \& {}_y\backslash Y \arrow[d]\\
X_i \arrow[r]\arrow[d] \& X \arrow[r,"\unsure{p|r}"]\arrow[d,"f"'] \& Y\\
\substack{\{i\}\\[-1pt]\simeq e} \arrow[r] \& I \& {,}
\end{tikzcd}}
où les trois carrés sont cartésiens. Comme $\,{}_y\backslash Y\unsure{-|\to}Y$ est $W$-propre, il en est de même de $\,{}_y\backslash X-X$ par chgt. de base, \unsure{donc} du composé $\,{}_y\backslash X-X\longrightarrow I$ par transitivité. Le diagramme montre d'autre part que la fibre de $\,{}_y\backslash X$ en $i\in\mathrm{Ob}\,I$ s'identifie à $\unsure{{}_y\backslash X_i}$.
